%=====================================================================
%  5. Results
%=====================================================================
\section{Results}
\label{sec:results}

\subsection{Pilot Study: Three-Case Probe}

Before committing to a full evaluation, we ran a three-case stratified probe
(one case per difficulty level, all on the same topic) to verify that the
pipeline could be measured end to end. The probe produced an informative
negative result.

\begin{table}[!t]
\centering
\caption{Pilot probe ($n=3$). Expected labels follow the public rule of
Section~\ref{ssec:diagnosis}; ``planned'' is the supervisor's output.}
\label{tab:pilot}
\begin{tabular}{@{}llcccc@{}}
\toprule
\textbf{Case} & \textbf{Profile} & \textbf{Expected} & \textbf{Planned} & \textbf{Match} & \textbf{Coverage} \\
\midrule
TC-002 & LP-02 & \textsc{basic}        & \textsc{basic}        & \checkmark & 5/6 \\
TC-001 & LP-01 & \textsc{intermediate} & \textsc{basic}        & $\times$   & 4/6 \\
TC-005 & LP-09 & \textsc{advanced}     & \textsc{intermediate} & $\times$   & 4/6 \\
\bottomrule
\end{tabular}
\end{table}

Adaptation accuracy was $1/3 = 33.3\%$, coverage $13/18 = 72.2\%$, and the
internally reported hallucination rate $0.00\%$ ($n=3$, too small to support
any reliability claim). Adaptation accuracy is far below the $85\%$ target,
and---critically---the errors are not random. They form the pattern shown in
Table~\ref{tab:pilot_confusion}: every case is either correct or planned
\emph{one level too low}. A single-direction offset of this kind indicates a
\emph{systematic calibration defect}, not sampling noise.

\begin{table}[!t]
\centering
\caption{Pilot difficulty confusion matrix (expected $\rightarrow$ planned).
The triangle structure reveals a systematic downward shift.}
\label{tab:pilot_confusion}
\begin{tabular}{@{}lccc@{}}
\toprule
\textbf{Expected $\downarrow$ / Planned $\rightarrow$} & \textsc{basic} & \textsc{inter.} & \textsc{adv.} \\
\midrule
\textsc{basic}        & 1 & 0 & 0 \\
\textsc{intermediate} & 1 & 0 & 0 \\
\textsc{advanced}     & 0 & 1 & 0 \\
\bottomrule
\end{tabular}
\end{table}

\subsection{Diagnosis of the Calibration Defect}

We traced the defect to three compounding causes, all confirmed by
inspection of the running system.

\emph{(i) The supervisor's instruction contained no thresholds.} Its prompt
asked the model to ``comprehensively consider'' the learner's theory scores
when planning difficulty. Comprehensiveness is not an executable
specification: the model was not told which subjects matter for the topic, how
to aggregate them, or where the boundaries between levels lie. Absent
anchors, a language model regresses toward the conservative choice, which is
exactly the observed behavior.

\emph{(ii) The evaluation request omitted the modifier terms.} The public
calibration rule is
$d = \mathrm{clip}(b + \delta_g + \delta_r + \delta_j)$, yet the request
payload carried only a subset of profile fields. The $\delta$ terms were never
transmitted. Even a model that knew the rule could not have applied it. This
is a defect in the \emph{evaluation harness}, not in the model.

\emph{(iii) The system had no topic-to-subject mapping.} Nothing in the
codebase related teaching topics to their prerequisite subjects, so the
relevant subset of the learner's scores was undetermined.

\subsection{Calibration Instruction Ablation}
\label{ssec:calib_ablation}

The diagnosis suggests a testable intervention. We compared two supervisor
prompt conditions on the identical case set:

\begin{itemize}
\item[(a)] \textbf{Implicit} --- the original instruction, which merely asks
the model to consider the profile.
\item[(b)] \textbf{Explicit} --- the aggregation rule, the modifier terms, the
level boundaries, and the intermediate quantities ($b$, $d$) are computed and
injected into the prompt, with an escape hatch: the model may deviate from the
computed level only if it states concrete evidence, which is logged.
\end{itemize}

Condition (a) yields $33.3\%$ adaptation accuracy with the systematic
downward shift of Table~\ref{tab:pilot_confusion}. Condition (b) is expected
to remove the shift; the full-set measurement is reported in
Table~\ref{tab:calib}. \todo{Fill Table~\ref{tab:calib} after the
post-intervention run. Report both the accuracy change and, importantly,
whether the error direction is no longer one-sided.}

\begin{table}[!t]
\centering
\caption{Calibration instruction ablation.}
\label{tab:calib}
\begin{tabular}{@{}lccc@{}}
\toprule
\textbf{Condition} & \textbf{Acc.} & \textbf{Down-shifted errors} & \textbf{Over-shifted errors} \\
\midrule
(a) Implicit instruction & 33.3\% ($n$=3) & 2 & 0 \\
(b) Explicit rule + escape hatch & \todo{XX\%} & \todo{XX} & \todo{XX} \\
\bottomrule
\end{tabular}
\end{table}

We consider this ablation a contribution in its own right rather than a bug
report. It isolates a general failure mode of LLM-based decision modules in
educational systems: \textbf{when a decision must be reproducible, stating the
criteria is not optional.} The escape hatch matters too---pure rule lookup
would discard information the profile contains, while unconstrained judgment
is unauditable; the hybrid keeps both the accuracy and the
auditability.

\subsection{RQ1--RQ3: Main Comparison}

Table~\ref{tab:main} reports the primary comparison on the full case set.
\todo{Replace every ``XX'' with values produced by
\texttt{tests/run\_live\_parallel.py --all --concurrency 4}.}

\begin{table*}[!t]
\centering
\caption{Main results over \todo{61} end-to-end cases. Best value per column
in bold. $\uparrow$ higher is better, $\downarrow$ lower is better.
$\Delta$ is the difference against the full configuration.}
\label{tab:main}
\begin{tabular}{@{}lccccc@{}}
\toprule
\textbf{Config} & \textbf{Adaptation acc.\ $\uparrow$} & \textbf{Coverage $\uparrow$} &
\textbf{Halluc.\ rate $\downarrow$} & \textbf{Review overall $\uparrow$} & \textbf{$\Delta$ halluc.\ (pp)} \\
\midrule
B1 single LLM, no profile      & XX.X\% & XX.X\% & XX.X\% & X.X/10 & +XX.X \\
B2 single LLM + profile        & XX.X\% & XX.X\% & XX.X\% & X.X/10 & +XX.X \\
B3 multi-agent, no retrieval   & XX.X\% & XX.X\% & XX.X\% & X.X/10 & +XX.X \\
B4 multi-agent, no debate      & XX.X\% & XX.X\% & XX.X\% & X.X/10 & +XX.X \\
B5 self-refine instead         & XX.X\% & XX.X\% & XX.X\% & X.X/10 & +XX.X \\
B6 no review-gated retry       & XX.X\% & XX.X\% & XX.X\% & X.X/10 & +XX.X \\
\midrule
\textbf{Full (\sysname{})}     & \textbf{XX.X\%} & \textbf{XX.X\%} & \textbf{X.X\%} & \textbf{X.X/10} & --- \\
\bottomrule
\end{tabular}
\end{table*}

A methodological note on B4: with debate disabled the system no longer
\emph{self-reports} a hallucination rate. That rate must therefore be
measured by the same external judge used for every other condition. Reporting
a self-computed zero for an unverified configuration would be a measurement
artifact of exactly the kind this paper argues against.

\subsection{RQ4: Is the Internal Verifier Trustworthy?}

This question must be answered before any other number is believed, because
the verifier and the generator share a model family. Table~\ref{tab:agreement}
cross-tabulates the internal verdict against human raters and against external
model judges.

\begin{table}[!t]
\centering
\caption{Agreement on hallucination judgments (per claim).}
\label{tab:agreement}
\begin{tabular}{@{}lc@{}}
\toprule
\textbf{Comparison} & \textbf{Agreement} \\
\midrule
Internal vs.\ human  & XX.X\% ($\kappa$ = 0.XX) \\
External vs.\ human  & XX.X\% ($\kappa$ = 0.XX) \\
Internal vs.\ external & XX.X\% ($\kappa$ = 0.XX) \\
Human--human (rater pair) & XX.X\% ($\kappa$ = 0.XX) \\
\bottomrule
\end{tabular}
\end{table}

\todo{Interpretation to be written once numbers exist. If internal--human
agreement is materially lower than human--human agreement, the internal
verifier over- or under-reports, and the headline hallucination rate must be
restated using the human-anchored estimate. Do not suppress this case.}

\subsection{RQ5: Human Evaluation}

\begin{table}[!t]
\centering
\caption{Human ratings (5-point Likert, mean $\pm$ SD) on \todo{30} bundles.}
\label{tab:human}
\begin{tabular}{@{}lcc@{}}
\toprule
\textbf{Dimension} & \textbf{\sysname{}} & \textbf{Best baseline} \\
\midrule
Factual correctness      & X.XX $\pm$ X.XX & X.XX $\pm$ X.XX \\
Difficulty fit           & X.XX $\pm$ X.XX & X.XX $\pm$ X.XX \\
Instructional usefulness & X.XX $\pm$ X.XX & X.XX $\pm$ X.XX \\
\bottomrule
\end{tabular}
\end{table}

\subsection{Cost and Latency}

\begin{table}[!t]
\centering
\caption{Computational cost per case (mean).}
\label{tab:cost}
\begin{tabular}{@{}lc@{}}
\toprule
\textbf{Metric} & \textbf{Value} \\
\midrule
Wall-clock per case & \todo{XX.X s} \\
LLM calls per case  & \todo{XX} \\
Tokens per case     & \todo{XX,XXX} \\
Retry rate          & \todo{XX.X\%} \\
\bottomrule
\end{tabular}
\end{table}

The retry rate is worth noting: the bounded retry policy of
Algorithm~\ref{alg:verify} is designed to trigger rarely, so verification
overhead should stay far below a naive always-refine-twice design.
\todo{Confirm against measured retry rate.}
